%-------------------------
% Resume in LaTeX
% Author: Levon Khorasandzhian
% Track: Backend (.NET / Java)
%-------------------------

\documentclass[a4paper,11pt]{article}
\usepackage{../cvstyle}

\begin{document}

\begin{center}
    \textbf{\Huge \scshape Левон Хорасанджян} \\ \vspace{2pt}
    \small Backend-разработчик \resumeSeparator{} .NET / C\# / Java \\ \vspace{1pt}
    \small Москва, Россия \resumeSeparator{} Рассматриваю офис, гибрид и удаленный формат \\ \vspace{1pt}
    \resumeLink{mailto:lakhorasandzhyan@edu.hse.ru}{lakhorasandzhyan@edu.hse.ru} \resumeSeparator{}
    \resumeLink{https://t.me/yungflaeva}{t.me/yungflaeva} \resumeSeparator{}
    \resumeLink{tel:+79160171926}{+7 916 017-19-26} \resumeSeparator{}
    \resumeLink{https://github.com/lkhorasandzhian}{github.com/lkhorasandzhian}
\end{center}

\section{Образование}
  \resumeSubHeadingListStart
    \resumeSubheading
      {Высшая школа экономики}{Москва, Россия}
      {Бакалавриат Факультета Компьютерных Наук, программная инженерия}{Авг 2021 -- Июл 2025}
    \resumeSubSubheading
      {Магистратура Факультета Компьютерных Наук, инженерия данных}{Окт 2025 -- Июл 2027}
    \resumeItem{
      \textbf{Релевантные курсы (GPA: 8.5/10)}: Алгоритмы и структуры данных, Проектирование архитектуры программных систем, Программирование на C\#, Конструирование программного обеспечения (Java), Распределенные базы и хранилища данных, Нереляционные базы данных, Архитектура вычислительных систем, Разработка и анализ требований, Инструменты разработки ПО, Контроль качества ПО и тестирование
    }
  \resumeSubHeadingListEnd

\section{Проекты}
  \resumeSubHeadingListStart
    \resumeProjectHeading
      {\href{https://github.com/lkhorasandzhian/room-booking-service}{\underline{\textbf{Room Booking Service}}} $|$ \emph{C\#, ASP.NET Core 10, PostgreSQL, Docker, Swagger, JWT, CI, Testing}}{Апр 2026}
      \resumeItemListStart
        \resumeItem{Разработал REST API для бронирования переговорных комнат: управление комнатами, расписания доступности, генерация слотов, создание и отмена бронирований, обзор бронирований для администратора}
        \resumeItem{Реализовал JWT-аутентификацию и разграничение доступа по ролям для user и admin}
        \resumeItem{Спроектировал слой данных на EF Core и PostgreSQL с миграциями и Docker Compose}
        \resumeItem{Реализовал авто-генерацию 30-минутных слотов на основе расписаний комнат и идемпотентную логику отмены бронирований}
        \resumeItem{Покрыл ключевые бизнес-сценарии тестами с xUnit, FluentAssertions и PostgreSQL Testcontainers}
        \resumeItem{Настроил GitHub Actions CI pipeline для restore, build, запуска unit- и integration-тестов, а также нагрузочных тестов PostgreSQL-backed API}
        \resumeItem{Добавил Swagger-документацию, Makefile-автоматизацию, seed-скрипты БД и NBomber load tests с порогами latency и success rate}
      \resumeItemListEnd

    \resumeProjectHeading
      {\href{https://github.com/lkhorasandzhian/ylab-university}{\underline{\textbf{Product Catalog Service}}} $|$ \emph{Java, Spring, PostgreSQL, Docker, Swagger, AOP, Web MVC}}{Ноя 2025}
      \resumeItemListStart
        \resumeItem{Разработал REST API для marketplace-системы с аутентификацией, ролями и авторизацией}
        \resumeItem{Реализовал хранение данных в PostgreSQL с Docker-based контейнеризацией}
        \resumeItem{Документировал API с помощью Swagger и Javadoc}
        \resumeItem{Реализовал cross-cutting concerns (логирование, валидация, безопасность) с использованием Spring AOP}
        \resumeItem{Развил проект за 5 итераций: от консольного приложения и сервлетов до полноценного Spring MVC приложения}
      \resumeItemListEnd
  \resumeSubHeadingListEnd

\section{Технические навыки}
  \begin{itemize}[leftmargin=0.15in,label={}]
    \small{\item{
      \resumeSkill{Языки}{C\#, Java} \\
      \resumeSkill{Фреймворки}{ASP.NET Core, Web API, Entity Framework Core, LINQ, Spring Core, Web MVC, REST, Security, AOP} \\
      \resumeSkill{Базы данных}{PostgreSQL, H2} \\
      \resumeSkill{Инструменты}{Git, Docker, Docker Compose, GitHub Actions, Linux, Postman, Swagger, Javadoc, Make} \\
      \resumeSkill{Тестирование}{xUnit, FluentAssertions, Testcontainers, NBomber, JUnit 5, Mockito}
    }}
  \end{itemize}

\section{Языки}
  \begin{itemize}[leftmargin=0.15in,label={}]
    \small{\item{
      \resumeSkill{Английский}{B2} \\
      \resumeSkill{Русский}{Родной}
    }}
  \end{itemize}

\end{document}
